% --- Archon physics formalization source begin ---
% archon:physics
% archon:covers PhyXMiniProblems/problem_phyx_mini_0073.lean
% archon:source-report reports/phyx_mini/problem_phyx_mini_0073.source.json

\chapter{Physics problem phyx\_mini\_0073}
\label{ch:PhyXMiniProblems_problem_phyx_mini_0073}

\paragraph{Problem source.}
\#\# Physical scenario

\textbackslash{}text\{White light is incident onto a \} 30\textasciicircum{}\textbackslash{}circ \textbackslash{}text\{ prism at the \} 40\textasciicircum{}\textbackslash{}circ \textbackslash{}text\{ angle shown in figure. Violet light emerges perpendicular to the rear face of the prism. The index of refraction of violet light in this glass is \} 2.0\textbackslash{}\% \textbackslash{}text\{ larger than the index of refraction of red light.\}

\#\# Question

At what angle \textbackslash{}phi does red light emerge from the rear face?

\#\# Answer choices

- A: A: 3\textasciicircum{}\textbackslash{}circ

- B: B: 9\textasciicircum{}\textbackslash{}circ

- C: C: 1\textasciicircum{}\textbackslash{}circ

- D: D: 2\textasciicircum{}\textbackslash{}circ

\#\# Figure caption (auxiliary; use the image as primary evidence)

The image shows a triangular prism with a refraction diagram. 

- The prism is oriented with an apex angle labeled as \textbackslash{}(30\textasciicircum{}\textbackslash{}circ\textbackslash{}).
- An incident ray, labeled as "White light," enters the prism at an angle of \textbackslash{}(40\textasciicircum{}\textbackslash{}circ\textbackslash{}) with respect to the normal to the prism surface.
- Inside the prism, the light is refracted towards the base and exits the prism at another angle.
- The exiting light is split into two rays: one red and one blue, representing dispersion.
- The red ray exits at an angle labeled as \textbackslash{}(\textbackslash{}phi\textbackslash{}).
- The background is plain, and the triangle is shaded to represent the prism material.

\#\# Dataset metadata

Geometrical Optics; Physical Model Grounding Reasoning, Multi-Formula Reasoning

\paragraph{Recorded answer/context.}
C: C: 1\textasciicircum{}\textbackslash{}circ

\paragraph{Figure/image path.}
/root/proposal\_for\_physic/science-mango/phyx\_mini\_run/phyx\_data/test\_image/73.png

\paragraph{Formalization target.}
create a compiling Lean file with sorry bodies at `PhyXMiniProblems/problem\_phyx\_mini\_0073.lean`. The Lean declarations must preserve the physical quantities, dimensions or dimensional roles, figure labels, governing-law hypotheses, and final relation expressed by this problem.
Use Mathlib/Physlib names found through LeanExplore where available. If a physics API is missing, introduce faithful local abstractions rather than scalar placeholder aliases.

\begin{theorem}[Physics formalization target]
\label{thm:physics:phyx_mini_0073:target}
\lean{PhyXMiniProblems.ProblemPhyXMini0073.redLightEmergenceAngle_is_choiceC}
\uses{def:physics:phyx-mini-0073:phyxminiproblems-problemphyxmini0073-dispersiveprismsetup, def:physics:phyx-mini-0073:phyxminiproblems-problemphyxmini0073-hasphysicalopticalparameters, def:physics:phyx-mini-0073:phyxminiproblems-problemphyxmini0073-matchesproblemandfiguredata, def:physics:phyx-mini-0073:phyxminiproblems-problemphyxmini0073-satisfiesprismgeometry, def:physics:phyx-mini-0073:phyxminiproblems-problemphyxmini0073-obeyssnelllaw, def:physics:phyx-mini-0073:phyxminiproblems-problemphyxmini0073-phi, def:physics:phyx-mini-0073:phyxminiproblems-problemphyxmini0073-roundstoanswerchoice, def:physics:phyx-mini-0073:phyxminiproblems-problemphyxmini0073-isuniqueclosestanswerchoice}
The assigned autoformalize agent should translate this physics problem into Lean declarations in the covered file, with theorem and lemma proof bodies written as `by sorry`.
\end{theorem}
\begin{proof}
This is an autoformalization task, not a proof task. Produce statements that can later be proved by the physics prover without weakening the physical model.
\end{proof}

% --- Archon declaration topology begin ---
\section{Declaration topology}
\begin{definition}[radians Of Degrees]
  \label{def:physics:phyx-mini-0073:phyxminiproblems-problemphyxmini0073-radiansofdegrees}
  \lean{PhyXMiniProblems.ProblemPhyXMini0073.radiansOfDegrees}
  Convert a degree readout to its real radian representative
\end{definition}
\begin{proof}
  This is the declared model definition of radians Of Degrees.
\end{proof}
\begin{definition}[angle Of Degrees]
  \label{def:physics:phyx-mini-0073:phyxminiproblems-problemphyxmini0073-angleofdegrees}
  \lean{PhyXMiniProblems.ProblemPhyXMini0073.angleOfDegrees}
  \uses{def:physics:phyx-mini-0073:phyxminiproblems-problemphyxmini0073-radiansofdegrees}
  Convert a degree readout to Mathlib's type of angles modulo one turn
\end{definition}
\begin{proof}
  This is the declared model definition of angle Of Degrees.
\end{proof}
\begin{definition}[degrees Of Angle]
  \label{def:physics:phyx-mini-0073:phyxminiproblems-problemphyxmini0073-degreesofangle}
  \lean{PhyXMiniProblems.ProblemPhyXMini0073.degreesOfAngle}
  Principal degree readout of an angle, used to compare answer choices
\end{definition}
\begin{proof}
  This is the declared model definition of degrees Of Angle.
\end{proof}
\begin{definition}[Spectral Color]
  \label{def:physics:phyx-mini-0073:phyxminiproblems-problemphyxmini0073-spectralcolor}
  \lean{PhyXMiniProblems.ProblemPhyXMini0073.SpectralColor}
  The two spectral components singled out by the dispersion diagram
\end{definition}
\begin{proof}
  This is the declared model definition of Spectral Color.
\end{proof}
\begin{definition}[Optical Medium]
  \label{def:physics:phyx-mini-0073:phyxminiproblems-problemphyxmini0073-opticalmedium}
  \lean{PhyXMiniProblems.ProblemPhyXMini0073.OpticalMedium}
  The two homogeneous optical media traversed by each colored ray
\end{definition}
\begin{proof}
  This is the declared model definition of Optical Medium.
\end{proof}
\begin{definition}[Dispersive Prism Setup]
  \label{def:physics:phyx-mini-0073:phyxminiproblems-problemphyxmini0073-dispersiveprismsetup}
  \lean{PhyXMiniProblems.ProblemPhyXMini0073.DispersivePrismSetup}
  \uses{def:physics:phyx-mini-0073:phyxminiproblems-problemphyxmini0073-spectralcolor, def:physics:phyx-mini-0073:phyxminiproblems-problemphyxmini0073-opticalmedium}
  The physical quantities and named angles in the prism figure. The incident white ray has one common entrance incidence angle. Dispersion in the glass gives color-dependent internal and emerging angles. The field incidentRayToEntranceFace is the angle marked 40° between the incoming ray and the entrance *face*. All other path angles are measured from the indicated face normal. In particular, the red value of emergenceFromRearNormal is the figure's unknown φ, since the violet ray follows that normal
\end{definition}
\begin{proof}
  This is the declared model definition of Dispersive Prism Setup.
\end{proof}
\begin{definition}[Is Principal Optical Angle]
  \label{def:physics:phyx-mini-0073:phyxminiproblems-problemphyxmini0073-isprincipalopticalangle}
  \lean{PhyXMiniProblems.ProblemPhyXMini0073.IsPrincipalOpticalAngle}
  An angle lies on the principal geometrical-optics branch
\end{definition}
\begin{proof}
  This is the declared model definition of Is Principal Optical Angle.
\end{proof}
\begin{definition}[Has Physical Optical Parameters]
  \label{def:physics:phyx-mini-0073:phyxminiproblems-problemphyxmini0073-hasphysicalopticalparameters}
  \lean{PhyXMiniProblems.ProblemPhyXMini0073.HasPhysicalOpticalParameters}
  \uses{def:physics:phyx-mini-0073:phyxminiproblems-problemphyxmini0073-dispersiveprismsetup, def:physics:phyx-mini-0073:phyxminiproblems-problemphyxmini0073-isprincipalopticalangle}
  Positivity and branch information for the physical optical data. These conditions make the normal-referenced sine equations select physical angles
\end{definition}
\begin{proof}
  This is the declared model definition of Has Physical Optical Parameters.
\end{proof}
\begin{definition}[Matches Problem And Figure Data]
  \label{def:physics:phyx-mini-0073:phyxminiproblems-problemphyxmini0073-matchesproblemandfiguredata}
  \lean{PhyXMiniProblems.ProblemPhyXMini0073.MatchesProblemAndFigureData}
  \uses{def:physics:phyx-mini-0073:phyxminiproblems-problemphyxmini0073-angleofdegrees, def:physics:phyx-mini-0073:phyxminiproblems-problemphyxmini0073-dispersiveprismsetup}
  Numerical and figure-derived data supplied by the problem. The image makes the 40° mark face-referenced; its complementary normal-referenced incidence is imposed separately by prism geometry. No red emergence value occurs here
\end{definition}
\begin{proof}
  This is the declared model definition of Matches Problem And Figure Data.
\end{proof}
\begin{definition}[Satisfies Prism Geometry]
  \label{def:physics:phyx-mini-0073:phyxminiproblems-problemphyxmini0073-satisfiesprismgeometry}
  \lean{PhyXMiniProblems.ProblemPhyXMini0073.SatisfiesPrismGeometry}
  \uses{def:physics:phyx-mini-0073:phyxminiproblems-problemphyxmini0073-angleofdegrees, def:physics:phyx-mini-0073:phyxminiproblems-problemphyxmini0073-dispersiveprismsetup}
  The angle relations determined by the triangular figure. The face-marked entrance angle and normal-referenced incidence are complementary. On the depicted branch, the red internal ray is just above the rear-face normal, so its entrance-normal angle is the apex angle plus its rear incidence magnitude; the same equation specializes to zero rear incidence for violet
\end{definition}
\begin{proof}
  This is the declared model definition of Satisfies Prism Geometry.
\end{proof}
\begin{definition}[Satisfies Entrance Snell Law]
  \label{def:physics:phyx-mini-0073:phyxminiproblems-problemphyxmini0073-satisfiesentrancesnelllaw}
  \lean{PhyXMiniProblems.ProblemPhyXMini0073.SatisfiesEntranceSnellLaw}
  \uses{def:physics:phyx-mini-0073:phyxminiproblems-problemphyxmini0073-spectralcolor, def:physics:phyx-mini-0073:phyxminiproblems-problemphyxmini0073-dispersiveprismsetup}
  Snell's law n₁ sin θ₁ = n₂ sin θ₂ at the air--glass entrance for one spectral component of the incident white light
\end{definition}
\begin{proof}
  This is the declared model definition of Satisfies Entrance Snell Law.
\end{proof}
\begin{definition}[Satisfies Rear Face Snell Law]
  \label{def:physics:phyx-mini-0073:phyxminiproblems-problemphyxmini0073-satisfiesrearfacesnelllaw}
  \lean{PhyXMiniProblems.ProblemPhyXMini0073.SatisfiesRearFaceSnellLaw}
  \uses{def:physics:phyx-mini-0073:phyxminiproblems-problemphyxmini0073-spectralcolor, def:physics:phyx-mini-0073:phyxminiproblems-problemphyxmini0073-dispersiveprismsetup}
  Snell's law at the glass--air rear face for one spectral component
\end{definition}
\begin{proof}
  This is the declared model definition of Satisfies Rear Face Snell Law.
\end{proof}
\begin{definition}[Obeys Snell Law]
  \label{def:physics:phyx-mini-0073:phyxminiproblems-problemphyxmini0073-obeyssnelllaw}
  \lean{PhyXMiniProblems.ProblemPhyXMini0073.ObeysSnellLaw}
  \uses{def:physics:phyx-mini-0073:phyxminiproblems-problemphyxmini0073-dispersiveprismsetup, def:physics:phyx-mini-0073:phyxminiproblems-problemphyxmini0073-satisfiesentrancesnelllaw, def:physics:phyx-mini-0073:phyxminiproblems-problemphyxmini0073-satisfiesrearfacesnelllaw}
  Both red and violet obey Snell's law at both prism faces
\end{definition}
\begin{proof}
  This is the declared model definition of Obeys Snell Law.
\end{proof}
\begin{definition}[phi]
  \label{def:physics:phyx-mini-0073:phyxminiproblems-problemphyxmini0073-phi}
  \lean{PhyXMiniProblems.ProblemPhyXMini0073.phi}
  \uses{def:physics:phyx-mini-0073:phyxminiproblems-problemphyxmini0073-dispersiveprismsetup}
  The figure label φ: red light's emergence angle from the rear-face normal
\end{definition}
\begin{proof}
  This is the declared model definition of phi.
\end{proof}
\begin{definition}[Answer Choice]
  \label{def:physics:phyx-mini-0073:phyxminiproblems-problemphyxmini0073-answerchoice}
  \lean{PhyXMiniProblems.ProblemPhyXMini0073.AnswerChoice}
  The answer-choice labels printed with the problem
\end{definition}
\begin{proof}
  This is the declared model definition of Answer Choice.
\end{proof}
\begin{definition}[angle Degrees]
  \label{def:physics:phyx-mini-0073:phyxminiproblems-problemphyxmini0073-answerchoice-angledegrees}
  \lean{PhyXMiniProblems.ProblemPhyXMini0073.AnswerChoice.angleDegrees}
  \uses{def:physics:phyx-mini-0073:phyxminiproblems-problemphyxmini0073-answerchoice}
  The whole-degree angle printed next to each answer choice
\end{definition}
\begin{proof}
  This is the declared model definition of angle Degrees.
\end{proof}
\begin{definition}[recorded Answer Choice]
  \label{def:physics:phyx-mini-0073:phyxminiproblems-problemphyxmini0073-recordedanswerchoice}
  \lean{PhyXMiniProblems.ProblemPhyXMini0073.recordedAnswerChoice}
  \uses{def:physics:phyx-mini-0073:phyxminiproblems-problemphyxmini0073-answerchoice}
  Dataset metadata recording answer C; this is not a theorem premise
\end{definition}
\begin{proof}
  This is the declared model definition of recorded Answer Choice.
\end{proof}
\begin{definition}[Rounds To Answer Choice]
  \label{def:physics:phyx-mini-0073:phyxminiproblems-problemphyxmini0073-roundstoanswerchoice}
  \lean{PhyXMiniProblems.ProblemPhyXMini0073.RoundsToAnswerChoice}
  \uses{def:physics:phyx-mini-0073:phyxminiproblems-problemphyxmini0073-degreesofangle, def:physics:phyx-mini-0073:phyxminiproblems-problemphyxmini0073-answerchoice}
  A computed angle rounds to the displayed whole-degree choice
\end{definition}
\begin{proof}
  This is the declared model definition of Rounds To Answer Choice.
\end{proof}
\begin{definition}[answer Choice Error Radians]
  \label{def:physics:phyx-mini-0073:phyxminiproblems-problemphyxmini0073-answerchoiceerrorradians}
  \lean{PhyXMiniProblems.ProblemPhyXMini0073.answerChoiceErrorRadians}
  \uses{def:physics:phyx-mini-0073:phyxminiproblems-problemphyxmini0073-radiansofdegrees, def:physics:phyx-mini-0073:phyxminiproblems-problemphyxmini0073-dispersiveprismsetup, def:physics:phyx-mini-0073:phyxminiproblems-problemphyxmini0073-phi, def:physics:phyx-mini-0073:phyxminiproblems-problemphyxmini0073-answerchoice}
  Angular error between φ and one displayed answer, measured in radians
\end{definition}
\begin{proof}
  This is the declared model definition of answer Choice Error Radians.
\end{proof}
\begin{definition}[Is Unique Closest Answer Choice]
  \label{def:physics:phyx-mini-0073:phyxminiproblems-problemphyxmini0073-isuniqueclosestanswerchoice}
  \lean{PhyXMiniProblems.ProblemPhyXMini0073.IsUniqueClosestAnswerChoice}
  \uses{def:physics:phyx-mini-0073:phyxminiproblems-problemphyxmini0073-dispersiveprismsetup, def:physics:phyx-mini-0073:phyxminiproblems-problemphyxmini0073-answerchoice, def:physics:phyx-mini-0073:phyxminiproblems-problemphyxmini0073-answerchoiceerrorradians}
  A choice is strictly closer to φ than each of the other displayed angles
\end{definition}
\begin{proof}
  This is the declared model definition of Is Unique Closest Answer Choice.
\end{proof}
% --- Archon declaration topology end ---
% --- Archon physics formalization source end ---
